\documentclass{article}
\usepackage{amsmath,amssymb}
\usepackage{xurl}
\usepackage[hidelinks]{hyperref}
% Shared commands for notes in docs/paper-gaps/.

\DeclareUnicodeCharacter{2080}{\ifmmode{}_{0}\else\textsubscript{0}\fi}
\DeclareUnicodeCharacter{2081}{\ifmmode{}_{1}\else\textsubscript{1}\fi}
\DeclareUnicodeCharacter{2082}{\ifmmode{}_{2}\else\textsubscript{2}\fi}
\DeclareUnicodeCharacter{2099}{\ifmmode{}_{n}\else\textsubscript{n}\fi}

\newcommand{\Lean}{\textsc{Lean}}
\ExplSyntaxOn
\NewDocumentCommand{\leanid}{m}{
  \ifmmode
    \textsf{\detokenize{#1}}
  \else
    \group_begin:
    \urlstyle{sf}
    \tl_set:Nn \l_tmpa_tl {#1}
    \tl_replace_all:Nnn \l_tmpa_tl {\_} {_}
    \exp_args:NV \path \l_tmpa_tl
    \group_end:
  \fi
}
\ExplSyntaxOff
\providecommand{\tr}{\operatorname{tr}}
\newcommand{\tnleanIssueBase}{https://github.com/LionSR/TNLean/issues/}
\newcommand{\tnleanPrBase}{https://github.com/LionSR/TNLean/pull/}
\newcommand{\tnleanDocBase}{https://sirui-lu.com/TNLean/docs/}
\newcommand{\ghissue}[1]{\href{\tnleanIssueBase#1}{GitHub issue \##1}}
\newcommand{\ghpr}[1]{\href{\tnleanPrBase#1}{PR \##1}}
\newcommand{\doclink}[1]{\href{\tnleanDocBase#1.html}{\path{#1}}}

% Dirac notation and the identity operator, as in blueprint/src/macros/common.tex.
% \providecommand keeps note-local definitions authoritative.
\usepackage{dsfont}
\providecommand{\ket}[1]{|#1\rangle}
\providecommand{\bra}[1]{\langle #1|}
\providecommand{\braket}[2]{\langle #1 | #2 \rangle}
\providecommand{\ketbra}[2]{|#1\rangle\!\langle #2|}
\providecommand{\Id}{\mathds{1}}
\providecommand{\one}{\mathds{1}}
\providecommand{\C}{\mathbb{C}}
\providecommand{\MN}[1]{M_{#1}(\C)}
\providecommand{\MD}{M_D(\C)}

% Verdict marker: \gapnote{<kind>}{<status>}. Kind is the policy
% classification (clarification, local-correction, scope-restriction,
% unfaithful, false-source, open-gap); status is open, wip (elimination
% underway), resolved, or historical. Machine-read by the site index;
% prints a verdict line.
\newcommand{\gapnote}[2]{\par\noindent\textbf{Verdict:} #1 (#2).\par}

\title{Preparation in Logarithmic Depth: Chain Lengths Divisible by the Block Length}
\author{}
\date{2026-09-27; resolved 2026-09-28}
\begin{document}
\maketitle
\gapnote{scope-restriction}{resolved}

\section{What the source states}
Let $A$ be a normal tensor of physical dimension $d$ and bond dimension $D$, and
let $\ket{\phi_N}$ be the normalized translation-invariant state it generates on
a ring of $N$ sites. Equation~(1) of~\cite{Malz2024Preparation} states that
$\ket{\phi_N}$ is prepared with error $\varepsilon$ by a circuit of local gates
of depth
\begin{align}
    T
     & =O(\log(N/\varepsilon)).
    \label{eq:mswc24_depth_log}
\end{align}
The error of a unit vector $\ket\psi$ is
$\varepsilon(\psi,\phi_N)=1-\lvert\braket{\psi}{\phi_N}\rvert$. The argument
blocks $q$ sites, prepares the approximating state
$\ket{\widetilde\phi_N}=(\bigotimes_{i=1}^{N/q}U_i)\ket\Omega$ of eq.~(10) of
the source in depth $O(q)$, bounds its error by
$O((N/q)e^{-\gamma q/\xi})$ (Lemma~1), and concludes that
$q=O(\log(N/\varepsilon))$ suffices. Equation~(10) has $N/q$ blocks of $q$
sites, so the argument as written takes $q$ to divide $N$. The proof of the
lower bound in the Supplemental Material of the source, which uses the same
blocking, says instead that the blocks are ``all of the same size, $q_N$,
except for the last one, which may be larger''.

\section{The restriction formerly adopted}
The formal statement first asked that the block length divide the chain
length.\footnote{Formal declarations:
    \leanid{MPSPreparation.exists_isPreparedInDepth_approximationError_le} and
    \leanid{MPSPreparation.exists_isPreparedInDepth_le_log_of_dvd} in
    \doclink{TNLean/MPS/Preparation/DepthLogBound}.}
There is $C$, depending only on $d$ and $D$, such that for every normal tensor
$A$ with $D\geq1$ there are $a>0$ and $b\geq1$, depending only on $A$, with the
following property. For $0<\varepsilon\leq1$, every $N\geq1$ and every $q$
dividing $N$ with
$q\geq a\log(N/\varepsilon)+b$, a unit vector $\ket\psi$ with
$\varepsilon(\psi,\phi_N)\leq\varepsilon$ is prepared in depth at most $Cq$. If
moreover $N\geq2$ and $q\leq2(a\log(N/\varepsilon)+b)$, for instance
$q=\lceil a\log(N/\varepsilon)+b\rceil$ when it divides $N$, the depth is at most
$c\log(N/\varepsilon)$ with $c$ depending only on $A$. This is
(\ref{eq:mswc24_depth_log}) for those pairs $(N,\varepsilon)$ for which such a
$q$ divides $N$. For other chain lengths, such as prime $N$ larger than the
admissible block lengths, the statement says nothing.

\section{Elimination}
Allow the last block to be larger, as in the proof of the lower bound: $N=(M-1)q+q'$ with
$q\leq q'<2q$. Three steps generalize.
\begin{enumerate}
    \item The approximating state is the product of the isometries $V_q$ on the
          first $M-1$ blocks and $V_{q'}$ on the last block, applied to the same
          product of pairs $\ket\omega$.
    \item Its overlap with $\ket{\phi_N}$ is the trace of the product of $M-1$
          mixed transfer matrices of $P_q$ and one of $P_{q'}$ against the fixed
          point, and the telescoping bound of Lemma~1 applies to a product of
          matrices each within $O(e^{-\gamma q/\xi})$ of the same idempotent.
    \item The circuit applies the block unitary of length $q'$ to the last block;
          its depth is $O(q')=O(q)$.
\end{enumerate}
With these, (\ref{eq:mswc24_depth_log}) holds for every $N\geq2$ with
$N\geq q$, where $q=\lceil a\log(N/\varepsilon)+b\rceil$. Chains shorter
than this block length, $N<q$, have no block of length $q$ and are not covered
by the three steps; they need a separate argument. The source does not treat
them: its proof of the lower bound takes $N$ ``such that we have a large number
of blocks''.

\section{Resolution}
The restriction is removed by the unequal blocks of the proof of the lower
bound, together with an exact preparation of chains shorter than the block
length.\footnote{Formal declarations:
    \leanid{MPSPreparation.exists_isPreparedInDepth_le_log_of_mpvState_ne_zero}
    and \leanid{MPSPreparation.exists_isPreparedInDepth_le_log} in
    \doclink{TNLean/MPS/Preparation/LogDepthPreparation}; the three steps are
    \leanid{MPSTensor.blockIsometryState},
    \leanid{MPSTensor.exists_blockApproximationError_le_mul}, and
    \leanid{MPSPreparation.exists_isPreparedInDepth_blockIsometryState}; the
    short chains are
    \leanid{MPSPreparation.exists_isPreparedInDepth_normalizedMPVState} and
    \leanid{MPSPreparation.exists_isPreparedInDepth_of_norm_eq_one}.}
Let $q=\lceil a\log(N/\varepsilon)+b\rceil$.

\paragraph{Chains with $N\ge q$.} Write $N=(M-1)q+q'$ with $M=\lfloor N/q\rfloor$
and $q'=q+(N\bmod q)$, so $q\le q'<2q$. The three steps of the elimination
plan above hold for any lengths $\ell_1,\dots,\ell_M$ of the blocks. The
state $(\bigotimes_kV_{\ell_k})\ket\Omega$ is a unit vector when the blocked
tensors are injective. Its overlap with $\ket{\phi_N(A)}$ is
$\operatorname{Tr}(\tau_{\ell_1}\cdots\tau_{\ell_M})$, with $\tau_\ell$ the mixed transfer
matrix of $P_\ell$ against $P_\infty$. The telescoping identity
\begin{align}
    X_1\cdots X_M-E^M
     & =\sum_{k=1}^{M}X_1\cdots X_{k-1}(X_k-E)E^{M-k}
    \label{eq:mswc24_depth_telescoping_product}
\end{align}
bounds it for any matrices with $\lVert X_k-E\rVert\le\delta$ and $E^2=E$, by
$c((1+c\delta)^M-1)$ as for equal factors. With $\lVert\tau_\ell-\tau_\infty\rVert
=O(e^{-\gamma\ell/\xi})\le O(e^{-\gamma q/\xi})$ for $\ell\ge q$ and
$\lvert\lambda_2\rvert\le1$, the error is $O(Me^{-\gamma q/\xi})$, which is at
most $\varepsilon$ as for equal blocks. The circuit applies a staircase of depth
$O(\ell_k)$ on each block, in parallel; the depth is $O(q')=O(q)$.

\paragraph{Chains with $N<q$.} Then $N<a\log(N/\varepsilon)+b$, and it suffices
to prepare $\ket{\phi_N}$ exactly in depth $O(N)$. When the $N$-site blocked
tensor $\mathcal B_N=VP$ is injective and $N\ge3D$,
$\phi_N(A)(s)=\operatorname{Tr}\mathcal B_N^s=\sum_{l,r}\bra sV\ket{l,r}\operatorname{Tr} P^{(l,r)}$, so
$\ket{\phi_N}$ is the state of a single block with the pair
$\omega(r,l)\propto\operatorname{Tr} P^{(l,r)}$ in place of the pair of the fixed point, and
the same circuit prepares it in depth $O(N)$. This is the sequential
preparation of a matrix product state. Below the injectivity length plus $3D$,
which depends only on $A$, every unit vector is a unitary applied to
$\ket{0\cdots0}$ on a bounded number of sites, a circuit of bounded depth.
In both regimes the depth is at most a constant times
$a\log(N/\varepsilon)+b\le(a+b/\log2)\log(N/\varepsilon)$.

The normalization of $\ket{\phi_N(A)}$ presupposes $\ket{\phi_N(A)}\neq0$; the
main theorem assumes it, and it holds for every sufficiently large $N$; see
\path{docs/paper-gaps/mswc24_depth_upper_bound_nonzero_state.tex}.

\bibliographystyle{alpha}
\bibliography{references}
\end{document}
